\par\Needspace{4\baselineskip}
\begin{center}
\hypertarget{native-033-article015}{}\pdfbookmark[2]{15.}{article-015}15.
\end{center}
Since the mean anomaly $M$, as we have seen, is completely determined by $v$ and $\varphi$, just as $v$ is by $M$ and $\varphi$, it is clear that if all three quantities are regarded as variables at the same time, their differential variations must satisfy an equation of condition worth deriving. Differentiating \SourceLeaf{034}equation VII of Art. 8 first gives $\frac{\mathrm{d}E}{\sin E}=\frac{\mathrm{d}v}{\sin v}-\frac{\mathrm{d}\varphi}{\cos\varphi}$. Differentiating equation XII in the same way gives $\mathrm{d}M=(1-e\cos E)\mathrm{d}E-\sin E\cos\varphi\,\mathrm{d}\varphi$. Eliminating $\mathrm{d}E$ from these differential equations, we obtain
\[
\mathrm{d}M=\frac{\sin E(1-e\cos E)}{\sin v}\,\mathrm{d}v-\left(\sin E\cos\varphi+\frac{\sin E(1-e\cos E)}{\cos\varphi}\right)\mathrm{d}\varphi
\]
or, substituting for $\sin E$ and $1-e\cos E$ their values from equations VIII and III,
\[
\mathrm{d}M=\frac{r r}{a a\cos\varphi}\,\mathrm{d}v-\frac{r(r+p)\sin v}{a a\cos\varphi^2}\,\mathrm{d}\varphi
\]
or, finally, expressing both coefficients in terms of $v$ and $\varphi$ alone,
\[
\mathrm{d}M=\frac{\cos\varphi^3}{(1+e\cos v)^2}\,\mathrm{d}v-\frac{(2+e\cos v)\sin v\cos\varphi^2}{(1+e\cos v)^2}\,\mathrm{d}\varphi.
\]
Conversely, if we regard $v$ as a function of the quantities $M$ and $\varphi$, the equation takes the form
\[
\mathrm{d}v=\frac{a a\cos\varphi}{r r}\,\mathrm{d}M+\frac{(2+e\cos v)\sin v}{\cos\varphi}\,\mathrm{d}\varphi
\]
or, introducing $E$ in place of $v$\,\raisebox{-.35ex}{.}
\[
\mathrm{d}v=\frac{a a\cos\varphi}{r r}\,\mathrm{d}M+\frac{a a}{r r}(2-e\cos E-e e)\sin E\,\mathrm{d}\varphi.
\]

\par\Needspace{4\baselineskip}
\begin{center}
\hypertarget{native-034-article016}{}\pdfbookmark[2]{16.}{article-016}16.
\end{center}
To determine the radius vector $r$ completely from $v$ and $\varphi$ or from $M$ and $\varphi$, we also need $p$ or $a$. Its differential will therefore consist of three terms. By differentiating equation II of Art. 8, we obtain
\[
\frac{\mathrm{d}r}{r}=\frac{\mathrm{d}p}{p}+\frac{e\sin v}{1+e\cos v}\,\mathrm{d}v-\frac{\cos\varphi\cos v}{1+e\cos v}\,\mathrm{d}\varphi.
\]
Setting $\frac{\mathrm{d}p}{p}=\frac{\mathrm{d}a}{a}-2\operatorname{tang}\varphi\,\mathrm{d}\varphi$ here (as follows from differentiating equation I) and expressing $\mathrm{d}v$ in terms of $\mathrm{d}M$ and $\mathrm{d}\varphi$ according to the preceding article, we obtain after the requisite simplifications
\[
\frac{\mathrm{d}r}{r}=\frac{\mathrm{d}a}{a}+\frac{a}{r}\operatorname{tang}\varphi\sin v\,\mathrm{d}M-\frac{a}{r}\cos\varphi\cos v\,\mathrm{d}\varphi
\]
or
\[
\mathrm{d}r=\frac{r}{a}\,\mathrm{d}a+a\operatorname{tang}\varphi\sin v\,\mathrm{d}M-a\cos\varphi\cos v\,\mathrm{d}\varphi.
\]

These formulas, like those we derived in the preceding article, are based on the assumption that $v$, $\varphi$ and $M$, or rather $\mathrm{d}v$, $\mathrm{d}\varphi$ and $\mathrm{d}M$, are \SourceLeaf{035}expressed in parts of the radius. If we wish to express the variations of the angles $v$, $\varphi$ and $M$ in arcseconds, we must either divide those parts of the formulas that involve $\mathrm{d}v$, $\mathrm{d}\varphi$ or $\mathrm{d}M$ by $206264{,}8$, or multiply those containing $\mathrm{d}r$, $\mathrm{d}p$ or $\mathrm{d}a$ by the same number. The formulas of the preceding article, which are homogeneous in this respect, therefore require no change.
